\documentclass[a4paper, 10pt]{article}

\usepackage{float}
\usepackage{amsmath,amsfonts}
\usepackage{amssymb}
\usepackage{arydshln}
\usepackage{tikz}
\usepackage{graphicx}
\usepackage{listings,jvlisting}
\usepackage{url}
\usepackage{enumerate}
\usepackage{multicol}
\usepackage{geometry}
\usepackage{fontspec}
\usepackage{unicode-math}
\usepackage{luatexja-fontspec}
\usepackage{cancel}
\usepackage{indentfirst}
\usepackage{sectsty}
\usepackage{tcolorbox}
\tcbuselibrary{breakable}
\usepackage[
  backend = biber,
  citestyle = authoryear-comp, 
  bibstyle = authoryear, 
  date = year,
  dashed = false,
  eprint = false,
  url = false,
  isbn = false,
  maxcitenames = 2,
  maxbibnames = 10,
  uniquelist = minyear,
  uniquename = false,
  nohashothers = true,
  sorting=nyt, 
  sortcites=false,
  giveninits = true
]{biblatex}
\addbibresource{ref.bib}
\renewcommand*{\finalnamedelim}{\addspace\&\space}

\geometry{left=2cm, right=2cm, top=2cm, bottom=2cm}

\newtcolorbox{mybox}[1][]{%
    breakable,
    enhanced,
    colframe=blue!75!black,
    colback=blue!10,
    coltitle=black,
    fonttitle=\bfseries,
    title=My Box,
    watermark color=black!10,
    watermark text=\ifnum\value{tcbcolframe}=1 Continued on next page\else Continued from previous page\fi,
    #1
}

\setmainjfont{Noto Serif CJK JP}
\setmainfont{STIX-Two-Text}
%\setmainfont{CMU-Serif}
\setsansfont{Helvetica-Neue}

\setmathfont{STIX-Two-Math}
%\setmathfont{Asana-Math}
%\setmathfont{Libertinus-Math}
%\setmathfont{DejaVu-Math-TeX-Gyre}
%\setmathfont{XITS-Math}
%\setmathfont{STIX}

%\sectionfont{\rmfamily\large\bfseries}
%\subsectionfont{\rmfamily\normalsize\bfseries}
\sectionfont{\sffamily\gtfamily\Large\bfseries}
\subsectionfont{\sffamily\gtfamily\Large\bfseries}
\subsubsectionfont{\sffamily\gtfamily\large\bfseries}

\lstset{
  frame={tb},
  numberstyle={\ttfamily\scriptsize},
  stepnumber=1,
  numbers=left,
  language=Python,
  basicstyle=\ttfamily\scriptsize,
  keywordstyle=\color{blue},
  stringstyle=\color{red},
  commentstyle=\color{green},
  morecomment=[l][\color{magenta}]{\#}
  }


\setlength{\parindent}{1em}

\title{\vspace{-2.0cm}A note on the implementation of generalized R/T coefficient method}
\author{Harusato Kimura}

\begin{document}

\maketitle

\section{Overview}

This note summarizes the theory behind the program in this repository.
It covers
(i) the generalized reflection/transmission (R/T) coefficient method \parencite{chen1993} for the phase velocities of Rayleigh waves in a stack of homogeneous, isotropic and elastic layers over a half-space,
(ii) the family of secular functions used for models with low-velocity layers \parencite{he2006,wu2016,yuan2024},
(iii) the medium response and the relative power of the modes, which are evaluated from the residue of the surface response, and
(iv) the phase velocities and attenuation coefficients in viscoelastic layers, which are obtained from the complex roots of the same secular functions.
The matrices below are written in the same form as in \texttt{src/surface\_wave\_dispersion.cpp}.

\section{Model and notation}

The model consists of $N$ homogeneous layers ($i = 1, \dots, N$) over a homogeneous half-space ($i = N + 1$).
The $z$-axis points downward and the free surface is at $z = z^{(0)} = 0$.
The $i$th layer occupies $z^{(i - 1)} < z < z^{(i)}$, and its thickness is $h^{(i)} = z^{(i)} - z^{(i - 1)}$.
The P-wave velocity, S-wave velocity and density of the $i$th layer are $\alpha^{(i)}$, $\beta^{(i)}$ and $\rho^{(i)}$, respectively, and
\begin{align}
  \mu^{(i)} = \rho^{(i)} \left( \beta^{(i)} \right)^2, \qquad
  \lambda^{(i)} = \rho^{(i)} \left( \alpha^{(i)} \right)^2 - 2 \mu^{(i)}.
\end{align}
The superscript $(i)$ is omitted when a single layer is considered.
In the code, the arrays of the layer properties are zero-based, i.e., \texttt{var[i]} holds the value of the $(i + 1)$th layer, whereas \texttt{depth[i]} holds $z^{(i)}$.

\section{Motion-stress vector in a layer}

Consider a plane Rayleigh wave of angular frequency $\omega$ and horizontal wavenumber $k = \omega / c$ that propagates in the $+x$ direction,
\begin{align}
  u_x = \symrm{i}\, U(z)\, \symrm{e}^{\symrm{i}(\omega t - kx)}, \quad
  u_z = W(z)\, \symrm{e}^{\symrm{i}(\omega t - kx)}, \quad
  \tau_{zx} = \symrm{i}\, T(z)\, \symrm{e}^{\symrm{i}(\omega t - kx)}, \quad
  \tau_{zz} = S(z)\, \symrm{e}^{\symrm{i}(\omega t - kx)}.
\end{align}
With the factor $\symrm{i}$ in the horizontal components, $U$, $W$, $T$ and $S$ of a normal mode are real.
For real $k$, the results in this note do not depend on the choice of the sign of the time dependence.
Substitution into Hooke's law and the equation of motion gives a system of first-order differential equations for the motion-stress vector $\symbf{y} = (U, W, T, S)^{\symrm{T}}$,
\begin{align}
  \frac{\symrm{d} \symbf{y}}{\symrm{d} z} = A \symbf{y},
  \qquad
  A =
  \begin{pmatrix}
    0 & k & \dfrac{1}{\mu} & 0 \\[2mm]
    -\dfrac{\lambda k}{\lambda + 2\mu} & 0 & 0 & \dfrac{1}{\lambda + 2\mu} \\[2mm]
    \zeta k^2 - \rho \omega^2 & 0 & 0 & \dfrac{\lambda k}{\lambda + 2\mu} \\[2mm]
    0 & -\rho \omega^2 & -k & 0
  \end{pmatrix},
  \qquad
  \zeta = \frac{4 \mu (\lambda + \mu)}{\lambda + 2 \mu}.
\end{align}
$\symbf{y}$ is continuous across the interfaces.

In a homogeneous layer, the eigenvalues of $A$ are $\mp\gamma$ and $\mp\nu$, where
\begin{align}
  \gamma = \sqrt{k^2 - \frac{\omega^2}{\alpha^2}}, \qquad
  \nu = \sqrt{k^2 - \frac{\omega^2}{\beta^2}}, \qquad
  \symrm{Re}\,\gamma \geq 0, \quad \symrm{Re}\,\nu \geq 0.
\end{align}
The eigenvalues $-\gamma$ and $-\nu$ correspond to the down-going (decaying with depth) P and S waves, and $+\gamma$ and $+\nu$ to the up-going ones.
The general solution in the $i$th layer is therefore written as
\begin{align}
  \symbf{y}^{(i)}(z) =
  E^{(i)}
  \begin{pmatrix}
    \Lambda_d^{(i)}(z) & 0 \\
    0 & \Lambda_u^{(i)}(z)
  \end{pmatrix}
  \begin{pmatrix}
    \symbf{C}_d^{(i)} \\
    \symbf{C}_u^{(i)}
  \end{pmatrix},
  \qquad
  E^{(i)} =
  \begin{pmatrix}
    E_{11}^{(i)} & E_{12}^{(i)} \\
    E_{21}^{(i)} & E_{22}^{(i)}
  \end{pmatrix},
  \label{e:layer_solution}
\end{align}
where $\symbf{C}_d^{(i)}$ and $\symbf{C}_u^{(i)}$ are the two-component vectors of the amplitudes of the down-going and up-going (P, S) waves, and $E_{11}^{(i)}, \dots, E_{22}^{(i)}$ are $2 \times 2$ blocks.
The columns of $E$ are the eigenvectors of $A$,
\begin{align}
  E =
  \begin{pmatrix}
    -1 & \nu_1 & -1 & \nu_1 \\
    -\gamma_1 & 1 & \gamma_1 & -1 \\
    2 \mu \gamma & -\chi \mu & -2 \mu \gamma & \chi \mu \\
    \chi \mu & -2 \mu \nu & \chi \mu & -2 \mu \nu
  \end{pmatrix},
  \qquad
  \gamma_1 = \frac{\gamma}{k}, \quad
  \nu_1 = \frac{\nu}{k}, \quad
  \chi = 2k - \frac{\omega^2}{\beta^2 k},
\end{align}
and its inverse is
\begin{align}
  E^{-1} =
  \frac{k}{2 \rho \omega^2}
  \begin{pmatrix}
    -2 k \mu & \mu \chi / \gamma_1 & 1 / \gamma_1 & -1 \\
    -\mu \chi / \nu_1 & 2 k \mu & 1 & -1 / \nu_1 \\
    -2 k \mu & -\mu \chi / \gamma_1 & -1 / \gamma_1 & -1 \\
    -\mu \chi / \nu_1 & -2 k \mu & -1 & -1 / \nu_1
  \end{pmatrix},
\end{align}
where $\mu (4k - 2\chi) = 2 \rho \omega^2 / k$ was used for the prefactor.

\section{Generalized R/T coefficients}

\subsection{Phase factors}

The diagonal matrices in \eqref{e:layer_solution} are defined as
\begin{align}
  \Lambda_d^{(i)}(z) &= \symrm{diag}\left\{
  \symrm{e}^{-\gamma^{(i)} (z - z^{(i - 1)})},
  \symrm{e}^{-\nu^{(i)} (z - z^{(i - 1)})}
  \right\} 
  \\
  \Lambda_u^{(i)}(z) &= \symrm{diag}\left\{ 
  \symrm{e}^{-\gamma^{(i)}(z^{(i)} - z)},
  \symrm{e}^{-\nu^{(i)}(z^{(i)} - z)}
  \right\}
\end{align}

$\Lambda_d^{(i)}(z^{(i)})$ and $\Lambda_u^{(i + 1)}(z^{(i)})$ are expressed as
\begin{align}
  \Lambda_d^{(i)}(z^{(i)}) 
  &= 
  \symrm{diag}\left\{ 
  \symrm{e}^{-\gamma^{(i)}(z^{(i)} - z^{(i - 1)})},
  \symrm{e}^{-\nu^{(i)}(z^{(i)} - z^{(i - 1)})}
  \right\}  
  = 
  \symrm{diag}\left\{ 
  \symrm{e}^{-\gamma^{(i)}h^{(i)}},
  \symrm{e}^{-\nu^{(i)}h^{(i)}}
  \right\}  
  \\
  \Lambda_u^{(i + 1)}(z^{(i)}) 
  &= 
  \symrm{diag}\left\{ 
  \symrm{e}^{-\gamma^{(i + 1)}(z^{(i + 1)} - z^{(i)})},
  \symrm{e}^{-\nu^{(i + 1)}(z^{(i + 1)} - z^{(i)})}
  \right\}
  =
  \symrm{diag}\left\{ 
  \symrm{e}^{-\gamma^{(i + 1)}h^{(i + 1)}},
  \symrm{e}^{-\nu^{(i + 1)}h^{(i + 1)}}
  \right\}.
\end{align}

Define $\Lambda_{\symrm{kernel}}^{(i)}$, for each layer as
\begin{align}
  \Lambda_{\symrm{kernel}}^{(i)} = 
  \symrm{diag}\left\{ 
  \symrm{e}^{-\gamma^{(i)}h^{(i)}},
  \symrm{e}^{-\nu^{(i)}h^{(i)}}
  \right\}.  
\end{align}

With these definitions, $\symbf{C}_d^{(i)}$ is the amplitude of the down-going waves at the top of the layer and $\symbf{C}_u^{(i)}$ is that of the up-going waves at its bottom.
Because $\symrm{Re}\,\gamma \geq 0$ and $\symrm{Re}\,\nu \geq 0$, no element of the phase factors grows exponentially, which is the reason why the method is stable at high frequencies.
For the half-space ($h^{(N + 1)} = \infty$), $\Lambda_{\symrm{kernel}}^{(N + 1)} = 0$.

\subsection{Downward recursion}

The continuity of $\symbf{y}$ at $z = z^{(i)}$ reads
\begin{align}
  E^{(i)}
  \begin{pmatrix}
    \Lambda_d^{(i)}(z^{(i)}) \symbf{C}_d^{(i)} \\
    \symbf{C}_u^{(i)}
  \end{pmatrix}
  =
  E^{(i + 1)}
  \begin{pmatrix}
    \symbf{C}_d^{(i + 1)} \\
    \Lambda_u^{(i + 1)}(z^{(i)}) \symbf{C}_u^{(i + 1)}
  \end{pmatrix}.
  \label{e:continuity}
\end{align}
The generalized transmission and reflection coefficients for the waves incident from above, $\hat{T}_d^{(i)}$ and $\hat{R}_{du}^{(i)}$, are defined by
\begin{align}
  \symbf{C}_d^{(i + 1)} = \hat{T}_d^{(i)} \symbf{C}_d^{(i)}, \qquad
  \symbf{C}_u^{(i)} = \hat{R}_{du}^{(i)} \symbf{C}_d^{(i)}.
  \label{e:def_Rdu}
\end{align}
They are ``generalized'' in the sense that they include all the reverberations in the structure below.
There is no up-going wave in the half-space, $\symbf{C}_u^{(N + 1)} = 0$, so that
\begin{align}
  \hat{R}_{du}^{(N + 1)} = 0.
\end{align}
Substituting \eqref{e:def_Rdu} into \eqref{e:continuity} and writing
\begin{align}
  W^{(i)} = \left( E^{(i)} \right)^{-1} E^{(i + 1)} =
  \begin{pmatrix}
    W_{11}^{(i)} & W_{12}^{(i)} \\
    W_{21}^{(i)} & W_{22}^{(i)}
  \end{pmatrix},
\end{align}
we obtain the recursion from the bottom to the top ($i = N, N - 1, \dots, 1$),
\begin{align}
  \hat{T}_d^{(i)} &=
  \left[ W_{11}^{(i)} + W_{12}^{(i)} \Lambda_u^{(i + 1)}(z^{(i)}) \hat{R}_{du}^{(i + 1)} \right]^{-1}
  \Lambda_d^{(i)}(z^{(i)}),
  \\
  \hat{R}_{du}^{(i)} &=
  \left[ W_{21}^{(i)} + W_{22}^{(i)} \Lambda_u^{(i + 1)}(z^{(i)}) \hat{R}_{du}^{(i + 1)} \right]
  \hat{T}_d^{(i)}.
\end{align}

\subsection{Upward recursion}

In the same manner, the coefficients for the waves incident from below, $\hat{T}_u^{(i)}$ and $\hat{R}_{ud}^{(i)}$, are defined by
\begin{align}
  \symbf{C}_u^{(i)} = \hat{T}_u^{(i)} \symbf{C}_u^{(i + 1)}, \qquad
  \symbf{C}_d^{(i + 1)} = \hat{R}_{ud}^{(i)} \symbf{C}_u^{(i + 1)}.
  \label{e:def_Rud}
\end{align}
The traction vanishes at the free surface, i.e., the lower two rows of \eqref{e:layer_solution} vanish at $z = z^{(0)}$ for $i = 1$,
\begin{align}
  E_{21}^{(1)} \symbf{C}_d^{(1)} + E_{22}^{(1)} \Lambda_u^{(1)}(z^{(0)}) \symbf{C}_u^{(1)} = 0,
  \label{e:free_surface}
\end{align}
which gives the starting value
\begin{align}
  \hat{R}_{ud}^{(0)} = -\left( E_{21}^{(1)} \right)^{-1} E_{22}^{(1)} \Lambda_u^{(1)}(z^{(0)}).
\end{align}
With
\begin{align}
  O^{(i)} = \left( E^{(i + 1)} \right)^{-1} E^{(i)} =
  \begin{pmatrix}
    O_{11}^{(i)} & O_{12}^{(i)} \\
    O_{21}^{(i)} & O_{22}^{(i)}
  \end{pmatrix},
\end{align}
\eqref{e:continuity} and \eqref{e:def_Rud} give the recursion from the top to the bottom ($i = 1, 2, \dots$),
\begin{align}
  \hat{T}_u^{(i)} &=
  \left[ O_{21}^{(i)} \Lambda_d^{(i)}(z^{(i)}) \hat{R}_{ud}^{(i - 1)} + O_{22}^{(i)} \right]^{-1}
  \Lambda_u^{(i + 1)}(z^{(i)}),
  \\
  \hat{R}_{ud}^{(i)} &=
  \left[ O_{11}^{(i)} \Lambda_d^{(i)}(z^{(i)}) \hat{R}_{ud}^{(i - 1)} + O_{12}^{(i)} \right]
  \hat{T}_u^{(i)}.
\end{align}

\section{Secular functions}

\subsection{Secular function associated with the free surface}

Substituting $\symbf{C}_u^{(1)} = \hat{R}_{du}^{(1)} \symbf{C}_d^{(1)}$ into \eqref{e:free_surface}, a non-trivial $\symbf{C}_d^{(1)}$ exists only if
\begin{align}
  F_0(\omega, k) =
  \det \left[ E_{21}^{(1)} + E_{22}^{(1)} \Lambda_u^{(1)}(z^{(0)}) \hat{R}_{du}^{(1)} \right] = 0.
  \label{e:F0}
\end{align}
This is the dispersion equation of Rayleigh waves, and only the downward recursion is required to evaluate it.

\subsection{Secular functions associated with low-velocity layers}

A mode trapped in a low-velocity layer (LVL) has an almost vanishing amplitude at the free surface at high frequencies, and its root is then hardly detected with $F_0$.
Combining \eqref{e:def_Rdu} and \eqref{e:def_Rud} in the $l$th layer, $\symbf{C}_u^{(l)} = \hat{R}_{du}^{(l)} \symbf{C}_d^{(l)}$ and $\symbf{C}_d^{(l)} = \hat{R}_{ud}^{(l - 1)} \symbf{C}_u^{(l)}$, gives an equivalent dispersion equation that is evaluated inside that layer \parencite{chen1993,he2006,wu2016},
\begin{align}
  F_l(\omega, k) =
  \det \left[ I - \hat{R}_{ud}^{(l - 1)} \hat{R}_{du}^{(l)} \right] = 0.
\end{align}
In the code, the $i$th layer ($2 \leq i \leq N$) is regarded as an LVL if $\beta^{(i)} < \beta^{(i - 1)}$ and $\beta^{(i)} < \beta^{(i + 1)}$, and $F_0$ and $F_l$ of all the LVLs are used together.
The upward recursion is carried out only down to the deepest LVL, and it is skipped when the model has no LVL.

\subsection{Root search}

The roots are searched for in terms of the phase velocity $c = \omega / k$ at each frequency.
\begin{enumerate}
  \item The interval from $0.8 \min_i \beta^{(i)}$ to $\beta^{(N + 1)}$ is divided into $\lfloor (M + 1)\, n_{\symrm{grid}} \rfloor$ points, where $M$ is the number of modes predicted after \textcite{fan2007},
  \begin{align}
    M = \frac{\omega}{\pi c} \sum_{i = 1}^{N - 1}
    \symrm{Re} \left[
      \sqrt{ \left( \frac{c}{\alpha^{(i)}} \right)^2 - 1 }
      +
      \sqrt{ \left( \frac{c}{\beta^{(i)}} \right)^2 - 1 }
    \right] h^{(i)},
    \qquad c = \beta^{(N + 1)}.
  \end{align}
  Additional points are placed at $c_{\symrm{prev}} + \{0, \pm 5 \times 10^{-5}, \pm 5 \times 10^{-4}, \dots, \pm 5\}$ around every root $c_{\symrm{prev}}$ found at the previous frequency.
  \item The real and imaginary parts of $F_0$ and $F_l$ are evaluated at the grid points. If one of them changes its sign between two adjacent points, that component is bisected.
  \item The bisection is terminated when the width of the interval is smaller than $\epsilon_c = 10^{-3}$ (in the unit of the velocities of the input), and the candidate is accepted only if the root mean square of the real and imaginary parts of that secular function is sufficiently small there. A sign change that is caused by a singularity instead of a zero is rejected in this way.
  \item A candidate that coincides with $\alpha^{(i)}$ or $\beta^{(i)}$ of a layer is rejected, because $E^{(i)}$ is singular there.
\end{enumerate}
The value of the fundamental mode at zero frequency is linearly extrapolated from the first two frequencies.

\section{Medium response and relative power}

\subsection{Surface response}

Suppose that a traction $\symbf{\tau}_0 = (T(z^{(0)}), S(z^{(0)}))^{\symrm{T}}$ is applied on the free surface.
$\symbf{C}_u^{(1)} = \hat{R}_{du}^{(1)} \symbf{C}_d^{(1)}$ still holds because it follows from the conditions below the surface, so that the displacement $\symbf{u}_0 = (U(z^{(0)}), W(z^{(0)}))^{\symrm{T}}$ and the traction at the surface are
\begin{align}
  \symbf{u}_0 = P\, \symbf{C}_d^{(1)}, \qquad
  \symbf{\tau}_0 = Q\, \symbf{C}_d^{(1)},
\end{align}
with
\begin{align}
  P = E_{11}^{(1)} + E_{12}^{(1)} \Lambda_u^{(1)}(z^{(0)}) \hat{R}_{du}^{(1)}, \qquad
  Q = E_{21}^{(1)} + E_{22}^{(1)} \Lambda_u^{(1)}(z^{(0)}) \hat{R}_{du}^{(1)}.
\end{align}
The surface response $Y$ is defined as the matrix that relates them,
\begin{align}
  \symbf{u}_0 = Y \symbf{\tau}_0, \qquad
  Y(\omega, k) = P Q^{-1} =
  \begin{pmatrix}
    Y_{xx} & Y_{xz} \\
    Y_{zx} & Y_{zz}
  \end{pmatrix}.
\end{align}
Since $F_0 = \det Q$ from \eqref{e:F0}, $Y$ has a simple pole at every root $k = k_n$ of the dispersion equation.

\subsection{Medium response as a residue}

The medium response of the $n$th mode for a vertical point force and the vertical displacement at the free surface is defined as
\begin{align}
  A_n(\omega) = k_n \mathop{\symrm{Res}}_{k = k_n} Y_{zz}(\omega, k).
  \label{e:medium_response}
\end{align}
Let $U_n(z)$ and $W_n(z)$ be the eigenfunctions of the mode, $c_n$ and $U_{g,n}$ its phase and group velocities, and
\begin{align}
  I_{1,n} = \frac{1}{2} \int_0^\infty \rho \left( U_n^2 + W_n^2 \right) \symrm{d}z
\end{align}
its energy integral \parencite[eq.~7.74 of][]{aki2002}.
The residue of $Y$ is a matrix of rank one,
\begin{align}
  \mathop{\symrm{Res}}_{k = k_n} Y =
  -\frac{1}{4 k_n c_n U_{g,n} I_{1,n}}
  \begin{pmatrix}
    U_n(0)^2 & U_n(0) W_n(0) \\
    U_n(0) W_n(0) & W_n(0)^2
  \end{pmatrix},
  \qquad
  A_n = -\frac{W_n(0)^2}{4 c_n U_{g,n} I_{1,n}}.
  \label{e:residue_energy}
\end{align}
The negative sign appears because a force $f_z$ applied on the surface in the $+z$ direction corresponds to $\tau_{zz} = -f_z$ there.
\eqref{e:residue_energy} was confirmed numerically: for a model with one layer over a half-space, the two sides of the expression for $A_n$ agreed with a relative difference smaller than $10^{-7}$ for the first three modes, where the eigenfunctions were integrated numerically and the group velocity was obtained by a finite difference of the dispersion curve.

$|A_n|$ is the factor that appears in the surface-wave term of the Green function \parencite[eqs.~7.144 and 7.147 of][]{aki2002}.
The amplitude of the vertical displacement of the $n$th mode at a horizontal distance $r$ due to a vertical point force $F_z$ at the surface is
\begin{align}
  \left| u_z \right| =
  \left| F_z \right| \frac{W_n(0)^2}{8 c_n U_{g,n} I_{1,n}} \sqrt{\frac{2}{\pi k_n r}}
  =
  \left| F_z \right| \frac{\left| A_n \right|}{2} \sqrt{\frac{2}{\pi k_n r}}.
  \label{e:far_field}
\end{align}
$A_n$ corresponds to the medium response $A_R = N_R / (\partial F_R / \partial k)$ of \textcite{harkrider1964} in the Thomson--Haskell formulation \parencite{haskell1953}.
The ratio of the off-diagonal to the diagonal elements of the residue also gives the ellipticity, $U_n(0) / W_n(0) = \symrm{Res}\, Y_{xz} / \symrm{Res}\, Y_{zz}$.

\subsection{Evaluation by contour integration}

The residue in \eqref{e:medium_response} is evaluated by the contour integral along a circle of radius $r_c$ centred at the root, with the trapezoidal rule of $M_c$ points,
\begin{align}
  \mathop{\symrm{Res}}_{k = k_n} Y_{zz}
  = \frac{1}{2 \pi \symrm{i}} \oint Y_{zz}(k)\, \symrm{d}k
  \approx \frac{1}{M_c} \sum_{m = 0}^{M_c - 1}
  r_c\, \symrm{e}^{\symrm{i} \theta_m}\, Y_{zz}\!\left( k_n + r_c\, \symrm{e}^{\symrm{i} \theta_m} \right),
  \qquad
  \theta_m = \frac{2 \pi (m + 1/2)}{M_c}.
\end{align}
Only the downward recursion with a complex wavenumber is required.
This approach has the following properties.
\begin{itemize}
  \item $Y$ does not change when the sign of $\gamma^{(i)}$ or $\nu^{(i)}$ of a finite layer is flipped, whereas $F_0$ does. $Y$ is hence analytic around a root on the real axis except at the other roots and at the branch points of the half-space, and the trapezoidal rule converges geometrically with $M_c$.
  \item The result does not depend on a step size, and it is insensitive to the error of the root as long as the root is inside the circle. A finite difference along the real axis has neither property; in a test, it failed for the modes trapped in an LVL, whose residues are many orders of magnitude smaller than the regular part of $Y$.
\end{itemize}
In the code, $M_c = 32$ and
\begin{align}
  r_c = \min \left( 2 \times 10^{-4}\, k_n,\ \frac{1}{4} \min_{l \neq n} \left| k_l - k_n \right| \right),
\end{align}
so that the circle does not enclose the other roots.
The root is known only within $\epsilon_c / 2$ in phase velocity, i.e., within $\delta k = k_n \epsilon_c / (2 c_n)$ in wavenumber.
If $r_c \leq 4\, \delta k$, the mode is regarded as unresolved and \texttt{nan} is written.

\subsection{Relative power}

The program outputs $|A_n|$ of every mode.
Its unit follows from those of the input; for lengths in m, velocities in m/s and densities in $\symrm{t/m^3}$, $\mu$ is in kPa and $|A_n|$ is in $1 / (\symrm{kPa\,m})$, so that \eqref{e:far_field} gives the displacement in m for a force in kN.
From \eqref{e:far_field}, the power of the $n$th mode at a given distance is proportional to $|A_n|^2 / k_n = |A_n|^2 c_n / \omega$.
The program also outputs the power relative to that of the mode with the lowest phase velocity at each frequency ($n = 0$),
\begin{align}
  \frac{P_n}{P_0} = \frac{\left| A_n \right|^2 c_n}{\left| A_0 \right|^2 c_0},
\end{align}
which is the same quantity as that of the Thomson--Haskell implementation.
For two test models (one layer over a half-space, and three layers including an LVL over a half-space), the two implementations agreed up to $25\,\symrm{Hz}$ with a median relative difference of about $10^{-7}$.
The ratio is not meaningful if the $0$th mode is trapped in an LVL and has a vanishing power at the surface.

\section{Viscoelastic medium}

If $Q_P^{(i)}$ and $Q_S^{(i)}$ are given in the input, the attenuation coefficients of the modes are computed instead of the medium responses and the relative powers.

\subsection{Complex velocities}

By the correspondence principle, all the expressions above hold for a viscoelastic medium if $\alpha^{(i)}$, $\beta^{(i)}$ and $\mu^{(i)}$ are replaced with complex values.
They are given by the constant-$Q$ model of \textcite{kjartansson1979}, in which a complex modulus $M$ depends on the frequency as
\begin{align}
  M(\omega) = M_0 \left( \frac{\symrm{i} \omega}{\omega_0} \right)^{2 \eta},
  \qquad
  \eta = \frac{1}{\pi} \arctan \frac{1}{Q}.
\end{align}
Let $c_0$ be the phase velocity at the reference frequency $\omega_0$, which is the velocity written in the input.
The complex velocity $\tilde{c} = \sqrt{M / \rho}$ at $\omega_0$ is then
\begin{align}
  \tilde{c}(\omega_0)
  = c_0 \cos \frac{\pi \eta}{2}\, \symrm{e}^{\symrm{i} \pi \eta / 2}
  = c_0\, \frac{1 + s + \symrm{i} / Q}{2 s},
  \qquad
  s = \sqrt{1 + Q^{-2}},
\end{align}
and both the phase velocity and the complex velocity are proportional to $(\omega / \omega_0)^{\eta}$.
This is applied to the P wave with $Q_P$ and to the S wave with $Q_S$ independently, and $\mu = \rho \tilde{\beta}^2$.
The imaginary part of $\tilde{c}$ is positive because the time dependence is $\symrm{e}^{\symrm{i} \omega t}$; unlike the elastic case, the sign of the time dependence matters here.
A layer with $Q \leq 0$ is treated as elastic.

In the code, $\tilde{c}(\omega_0)$ is used at all the frequencies unless the reference frequency is given, i.e., the complex velocities do not depend on the frequency.
If it is given, $\tilde{c}(\omega) = \tilde{c}(\omega_0) (\omega / \omega_0)^{\eta}$ is used, and the phase velocities $c_0 (\omega / \omega_0)^{\eta}$ are used as the velocities of the elastic medium at each frequency.
Since $\eta$ differs from layer to layer, the LVLs are then judged at each frequency.

\subsection{Complex roots}

A mode of a viscoelastic medium is a root of $F_0$ or $F_l$ at a complex wavenumber
\begin{align}
  k = \frac{\omega}{c} - \symrm{i} a,
\end{align}
where $c$ is the phase velocity and $a$ is the attenuation coefficient: the amplitude decays as $\symrm{e}^{-a x}$ in the direction of propagation.
The roots are searched for as follows at each frequency.
\begin{enumerate}
  \item The real roots $k_{\symrm{e}} = \omega / c_{\symrm{e}}$ of the elastic medium, i.e., of the medium with the real phase velocities, are obtained as in the root search above.
  \item Starting from $k = k_{\symrm{e}} - \symrm{i} a_{\symrm{prev}}$, where $a_{\symrm{prev}}$ is the attenuation coefficient of the same mode at the previous frequency ($0$ if it is not available), the Newton iteration $k \leftarrow k - F / (\partial F / \partial k)$ is applied to the secular function with which the elastic root was found. The secular functions are analytic in $k$ away from the branch cuts of $\nu$ and $\gamma$, so that the derivative is evaluated by a finite difference along the real axis with a step of $10^{-7} |k|$. A step is limited to $0.02 |k|$, and the iteration is terminated when the step becomes smaller than $10^{-13} |k|$.
  \item The root is accepted if $0.8 k_{\symrm{e}} < \symrm{Re}\, k < 1.2 k_{\symrm{e}}$ and $0 \leq a < 1$ (in the reciprocal of the unit of length of the input).
  \item A mode that converged to the same root as another mode (a relative difference smaller than $10^{-6}$), or that was not accepted, is solved again with every secular function from the initial guesses of the modes concerned. Among the roots that differ from those of all the other modes, the one with the smallest $|F / (\partial F / \partial k)| / |k|$ is taken. If there is none, the mode is written as nan at that frequency.
  \item The modes are numbered in ascending order of the phase velocity.
\end{enumerate}
The fourth step is needed where two modes cross, which they can do in a viscoelastic medium whereas they only approach each other in an elastic one.
In the cases examined so far, one of the two modes is trapped in an LVL there; its root almost cancels with a pole in $F_0$ and that of the other mode almost cancels with a pole in $F_l$, so that each secular function leads to only one of the two roots from both initial guesses.

If the lowest frequency is not zero, there is no previous frequency for the first one, and starting from $a = 0$ may lead to another root.
The attenuation is therefore switched on in ten steps there: the velocities are set to $c_0 (\tilde{c} / c_0)^{\xi}$ with $\xi = 0.1, 0.2, \dots, 0.9$, the root is followed by the Newton iteration, and the result is used as the initial guess.
The attenuation coefficient of the fundamental mode at zero frequency is linearly extrapolated from the first two frequencies.

\section{Rayleigh wave in a homogeneous half-space}

Secular function and characteristic equation for Rayleigh wave in homogeneous half-space is 
\begin{align}
  f(c) = 
  \left[ \left( \frac{c}{\alpha} \right)^2 - 1 \right]
  \left[ \left( \frac{c}{\beta} \right)^2 - 1 \right]
  - \left[ 1 - \frac{1}{2}\left( \frac{c}{\beta} \right)^2 \right]^4 = 0.
\end{align}

Its first and second derivatives are then expressed as 
\begin{align}
  \frac{\partial f}{\partial c}(c) &= 
  \frac{2}{\alpha}\left( \frac{c}{\alpha} \right)
  \left[ \left( \frac{c}{\beta} \right)^2 - 1 \right]
  +
  \frac{2}{\beta}\left( \frac{c}{\beta} \right)
  \left[ \left( \frac{c}{\alpha} \right)^2 - 1 \right]
  - 
  4 \left[\frac{1}{2}\left( \frac{c}{\beta} \right)^2 - 1 \right]^3 \frac{c}{\beta^2} 
  \\
  &= \frac{4}{\alpha^2 \beta^2}c^3 - 2\left( \frac{1}{\alpha^2} + \frac{1}{\beta^2} \right)c
  - 
  \frac{4c}{\beta^2} \left[\frac{1}{2}\left( \frac{c}{\beta} \right)^2 - 1 \right]^3 
\end{align}
and 
\begin{align}
  \frac{\partial^2 f}{\partial c^2}(c) 
  &= 
  \frac{12}{\alpha^2 \beta^2}c^2 - 2\left( \frac{1}{\alpha^2} + \frac{1}{\beta^2} \right)
  - 
  \frac{4}{\beta^2} \left[\frac{1}{2}\left( \frac{c}{\beta} \right)^2 - 1 \right]^3 
  - 
  \frac{12c^2}{\beta^4} \left[\frac{1}{2}\left( \frac{c}{\beta} \right)^2 - 1 \right]^2 
\end{align}


%\begin{figure}[H]
%  \begin{center}
%    \includegraphics[width=0.8\linewidth]{figs/case2.pdf}
%    \caption{Undamped time series obtained for model A and B with initial conditions of $d_0 = 1.0\,\symrm{m}$ and $v_0 = 10.0\,\symrm{ms^{-1}}$.}
%    \label{f:case1}
%  \end{center}
%\end{figure}

% -----------------------------------------------------------------------------
\printbibliography
% -----------------------------------------------------------------------------

\end{document}
